\documentclass{article}
%%%%%%%%%%%%%%%%%%%%%%%%%%%%%%%%%%%%%%%%%%%%%%%%%%%%%%%%%%%%%%%%%%%%%

\usepackage[upright]{fourier}% l'option permet d'avoir les majuscules droites dans les formules mathématiques
\usepackage{lmodern}
\usepackage[T1]{fontenc}
\usepackage[scale=.9]{geometry}
\usepackage{amssymb}
\usepackage{tikz}
\usepackage{setspace}

\usepackage{enumerate}
\renewcommand{\theenumi}{\bf\textsf{\arabic{enumi}}}
\renewcommand{\labelenumi}{\bf\textsf{\theenumi.}}
\renewcommand{\theenumii}{\bf\textsf{\alph{enumii}}}
\renewcommand{\labelenumii}{\bf\textsf{\theenumii.}}
\renewcommand{\theenumiii}{\bf\textsf{\roman{enumiii}}}
\renewcommand{\labelenumiii}{\bf\textsf{\theenumiii.}}

%les commandes suivantes évitent le message "too many math alphabets"...
\newcommand\hmmax{0}
\newcommand\bmmax{0}

%animaux
\usepackage{tikzlings}
%%%%%%%%%%%%%%%%%%%%%%%%%%%%%%%%%%%%%%%%%%%%%%%%%%%%%%%%%%%%%%%%%%%%%
%Titre

\title{QCM 1 - Fonction ln}
\author{\tikz[scale=1]\coati[sideward];
\includegraphics[width=0.2\linewidth]{Logo_augustin_thierry}
\tikz[scale=1]\coati[sideward,xscale=-1];
}
\date{\today} 

%%%%%%%%%%%%%%%%%%%%%%%%%%%%%%%%%%%%%%%%%%%%%%%%%%%%%%%%%%%%%%%%%%%%%
%Pour le QCM
\usepackage{hyperref}
\newcommand*{\mtcode}{123456}
%notice en ligne 
%http://www.tug.org/applications/hyperref/manual.html#x1-200006
%http://www.les-mathematiques.net/phorum/read.php?10,1653900,1654548

%%%%%%%%%%%%%%%%%%%%%%%%%%%%%%%%%%%%%%%%%%%%%%%%%%%%%%%%%%%%%%%%%%%%%
\begin{document}
%%%%%%%%%%%%%%%%%%%%%%%%%%%%%%%%%%%%%%%%%%%%%%%%%%%%%%%%%%%%%%%%%%%%%
\begin{spacing}{1.2}
%%%%%%%%%%%%%%%%%%%%%%%%%%%%%%%%%%%%%%%%%%%%%%%%%%%%%%%%%%%%%%%%%%%%

%%%%%%%%%%%%%%%%%%%%%%%%%%%%%%%%%%%%%%%
\maketitle

\begin{center}
\tikz[scale=0.6]\owl;
\tikz[scale=0.6]\owl;
\tikz[scale=0.6]\owl;
\tikz[scale=0.6]\owl;
\tikz[scale=0.6]\owl;
\tikz[scale=0.6]\owl;
\tikz[scale=0.6]\owl;
\tikz[scale=0.6]\owl;
\tikz[scale=0.6]\owl;
\tikz[scale=0.6]\owl;
\end{center}
%%%%%%%%%%%%%%%%%%%%%%%%%%%%%%%%%%%%%%%

\begin{Form}
\TextField[bordercolor=1 0 0,width=15em,charsize=0pt,name=box,bordercolor=black]{Prénom : }
\TextField[bordercolor=1 0 0,width=15em,charsize=0pt,name=box,bordercolor=black]{NOM : }


\begin{center}
\tikz[scale=0.6]\owl;
\tikz[scale=0.6]\owl;
\tikz[scale=0.6]\owl;
\tikz[scale=0.6]\owl;
\tikz[scale=0.6]\owl;
\tikz[scale=0.6]\owl;
\tikz[scale=0.6]\owl;
\tikz[scale=0.6]\owl;
\tikz[scale=0.6]\owl;
\tikz[scale=0.6]\owl;
\end{center}


\begin{enumerate}
\item La fonction ln est

\ChoiceMenu[radio,name=Q1,default=-0,align=0,radiosymbol=\ding{108},height=0.2cm,bordercolor=black]{}
{dérivable sur $]0;+\infty[$=1,
\hspace{0.2\linewidth}dérivable sur $\mathbb{R}$=0}

\item La fonction ln est

\ChoiceMenu[radio,name=Q2,default=-0,align=0,radiosymbol=\ding{108},height=0.2cm,bordercolor=black]{}
{dérivable sur $]0;+\infty[$=1,
\hspace{0.2\linewidth}dérivable sur $\mathbb{R}$=0}

\item La fonction ln est

\ChoiceMenu[radio,name=Q3,default=-0,align=0,radiosymbol=\ding{108},height=0.2cm,bordercolor=black]{}
{dérivable sur $]0;+\infty[$=1,
\hspace{0.2\linewidth}dérivable sur $\mathbb{R}$=0}

\end{enumerate}

\begin{center}
\tikz[scale=0.6]\owl;
\tikz[scale=0.6]\owl;
\tikz[scale=0.6]\owl;
\tikz[scale=0.6]\owl;
\tikz[scale=0.6]\owl;
\tikz[scale=0.6]\owl;
\tikz[scale=0.6]\owl;
\tikz[scale=0.6]\owl;
\tikz[scale=0.6]\owl;
\tikz[scale=0.6]\owl;
\end{center}

\TextField[password,bordercolor=1 0 0,width=10em,charsize=0pt,name=boxcode,bordercolor=black]{code : }
%
%
\PushButton[name=verificationcode,bordercolor=1 1 1,
onclick=
{
if (this.getField("boxcode").value==\mtcode)
{
this.getField("obtenirscore").value=(
this.getField("Q1").value
+this.getField("Q2").value
+this.getField("Q3").value
);
}
else
{
this.getField("obtenirscore").value="Code incorrect";
}
this.getField("boxcode").value="";
}
]
{\tikz\node[rounded corners, draw=black!80, fill=black!20]  {\bf score final :};}
\TextField[name=obtenirscore,bordercolor=1 1 1,width=10em,charsize=0pt,readonly=true]{}
\end{Form}

%%%%%%%%%%%%%%%%%%%%%%%%%%%%%%%%%%%%%%%%%%%%%%%%%%%%%%%%%%%%%%%%%%%%
\end{spacing}
%%%%%%%%%%%%%%%%%%%%%%%%%%%%%%%%%%%%%%%%%%%%%%%%%%%%%%%%%%%%%%%%%%%%
\end{document}